% GENERATED FILE. Edit canonical lessons, then run npm run generate:books.
% canonical-source-hash: fnv1a64:0aea0615e285c0fc
% canonical-lessons: GU-C364-sodh, GU-C364-avlokan, GU-C364-siddh, GU-C364-rahasya, GU-C364-ughadvu-uncover

\chapter{Discoveries}
\label{ch:gu-discoveries}
\hlchaptermodality{364}

\begin{chapteropening}
\textbf{By the end of this chapter:} \emph{I can say who found something out and how.}

Complete the last lesson of chapter 364: I can say who found something out and how.
\end{chapteropening}

\section[\texorpdfstring{\gu{શોધ} (śodh)}{શોધ (śodh)}]{\gu{શોધ} (śodh) — a discovery}

\label{lesson:GU-C364-sodh}



\begin{quote}
Before the new one: say the Gujarati for a count, then the Gujarati for to bring down.
\end{quote}

\subsection*{You'll want to know: \gu{શોધ}}
\textbf{\gu{શોધ}} — \emph{śodh} — \textquotedblleft{}a discovery\textquotedblright{}.

\begin{quote}
\gu{આ} \gu{દવા} \gu{એક} \gu{મોટી} \gu{શોધ} \gu{છે}.
\end{quote}

\begin{quote}
\emph{ā davā ek moṭī śodh che.}
\end{quote}

\begin{quote}
This medicine is a major discovery.
\end{quote}

A feminine noun.

\begin{cousinweb}
\textbf{\gu{શોધ}} (\emph{śodh}) is the noun of \textbf{\gu{શોધવું}} (\emph{śodhvũ}), \textquotedblleft{}to look for\textquotedblright{}, which CDIAL traces to Sanskrit \emph{śoddhum}, \textquotedblleft{}to purify\textquotedblright{} (CDIAL 12626): what a search turns up.
\end{cousinweb}

\begin{grammarlens}[title={What you've built}]
One more word for giving reasons and opinions.
\end{grammarlens}

\subsection*{Your turn}
\begin{itemize}
\raggedright
  \item \emph{Say it:} \emph{śodh}
  \item \emph{Say it:} \emph{śodh}, once more
  \item \emph{Say it:} \emph{ā davā ek moṭī śodh che.} — This medicine is a major discovery.
  \item \emph{Recall:} say the Gujarati for a count, then the Gujarati for to bring down, then say \emph{śodh} again
\end{itemize}

\begin{tcolorbox}[breakable,colback=teal!4,colframe=teal!35!black,title={Before you move on}]
What does \gu{શોધ} mean? (\textquotedblleft{}A discovery\textquotedblright{}.) Say it once more.
\end{tcolorbox}
\section[\texorpdfstring{\gu{અવલોકન} (avlokan)}{અવલોકન (avlokan)}]{\gu{અવલોકન} (avlokan) — an observation}

\label{lesson:GU-C364-avlokan}



\begin{quote}
Before the new one: say the Gujarati for to bring down, then the Gujarati for a discovery.
\end{quote}

\subsection*{You'll want to know: \gu{અવલોકન}}
\textbf{\gu{અવલોકન}} — \emph{avlokan} — \textquotedblleft{}an observation\textquotedblright{}.

\begin{quote}
\gu{મીરાએ} \gu{પોતાનાં} \gu{અવલોકનો} \gu{એક} \gu{ચોપડીમાં} \gu{નોંધ્યાં}.
\end{quote}

\begin{quote}
\emph{Mīrāe potānā̃ avlokno ek copḍīmā̃ nõdhyā̃.}
\end{quote}

\begin{quote}
Mira wrote her observations down in a notebook.
\end{quote}

A neuter noun.

\begin{cousinweb}
\textbf{\gu{અવલોકન}} (\emph{avlokan}) is a tatsama, borrowed from Sanskrit \emph{avalokana}, \textquotedblleft{}seeing, observing\textquotedblright{} (Monier-Williams).
\end{cousinweb}

\begin{grammarlens}[title={What you've built}]
One more word for giving reasons and opinions.
\end{grammarlens}

\subsection*{Your turn}
\begin{itemize}
\raggedright
  \item \emph{Say it:} \emph{avlokan}
  \item \emph{Say it:} \emph{avlokan}, once more
  \item \emph{Say it:} \emph{Mīrāe potānā̃ avlokno ek copḍīmā̃ nõdhyā̃.} — Mira wrote her observations down in a notebook.
  \item \emph{Recall:} say the Gujarati for to bring down, then the Gujarati for a discovery, then say \emph{avlokan} again
\end{itemize}

\begin{tcolorbox}[breakable,colback=teal!4,colframe=teal!35!black,title={Before you move on}]
What does \gu{અવલોકન} mean? (\textquotedblleft{}An observation\textquotedblright{}.) Say it once more.
\end{tcolorbox}
\section[\texorpdfstring{\gu{સિદ્ધ} (siddh)}{સિદ્ધ (siddh)}]{\gu{સિદ્ધ} (siddh) — proven}

\label{lesson:GU-C364-siddh}



\begin{quote}
Before the new one: say the Gujarati for a discovery, then the Gujarati for an observation.
\end{quote}

\subsection*{You'll want to know: \gu{સિદ્ધ}}
\textbf{\gu{સિદ્ધ}} — \emph{siddh} — \textquotedblleft{}proven\textquotedblright{}.

\begin{quote}
\gu{આ} \gu{પ્રયોગથી} \gu{તેની} \gu{વાત} \gu{સિદ્ધ} \gu{થઈ}.
\end{quote}

\begin{quote}
\emph{ā prayogthī tenī vāt siddh thaī.}
\end{quote}

\begin{quote}
This experiment proved what he said.
\end{quote}

\begin{cousinweb}
\textbf{\gu{સિદ્ધ}} (\emph{siddh}) is a tatsama, borrowed from Sanskrit \emph{siddha}, \textquotedblleft{}accomplished, established\textquotedblright{} (Monier-Williams).
\end{cousinweb}

\begin{grammarlens}[title={What you've built}]
One more word for giving reasons and opinions.
\end{grammarlens}

\subsection*{Your turn}
\begin{itemize}
\raggedright
  \item \emph{Say it:} \emph{siddh}
  \item \emph{Say it:} \emph{siddh}, once more
  \item \emph{Say it:} \emph{ā prayogthī tenī vāt siddh thaī.} — This experiment proved what he said.
  \item \emph{Recall:} say the Gujarati for a discovery, then the Gujarati for an observation, then say \emph{siddh} again
\end{itemize}

\begin{tcolorbox}[breakable,colback=teal!4,colframe=teal!35!black,title={Before you move on}]
What does \gu{સિદ્ધ} mean? (\textquotedblleft{}Proven\textquotedblright{}.) Say it once more.
\end{tcolorbox}
\section[\texorpdfstring{\gu{રહસ્ય} (rahasya)}{રહસ્ય (rahasya)}]{\gu{રહસ્ય} (rahasya) — a mystery; a secret}

\label{lesson:GU-C364-rahasya}



\begin{quote}
Before the new one: say the Gujarati for an observation, then the Gujarati for proven.
\end{quote}

\subsection*{You'll want to know: \gu{રહસ્ય}}
\textbf{\gu{રહસ્ય}} — \emph{rahasya} — \textquotedblleft{}a mystery; a secret\textquotedblright{}.

\begin{quote}
\gu{દાક્તરે} \gu{આ} \gu{રોગનું} \gu{રહસ્ય} \gu{શોધી} \gu{કાઢ્યું}.
\end{quote}

\begin{quote}
\emph{dāktare ā rognũ rahasya śodhī kāḍhyũ.}
\end{quote}

\begin{quote}
The doctor found out the secret of this illness.
\end{quote}

A neuter noun.

\begin{cousinweb}
\textbf{\gu{રહસ્ય}} (\emph{rahasya}) is a tatsama, borrowed from Sanskrit \emph{rahasya}, \textquotedblleft{}secret, mysterious; a secret, a mystery\textquotedblright{} (Monier-Williams).
\end{cousinweb}

\begin{grammarlens}[title={What you've built}]
One more word for giving reasons and opinions.
\end{grammarlens}

\subsection*{Your turn}
\begin{itemize}
\raggedright
  \item \emph{Say it:} \emph{rahasya}
  \item \emph{Say it:} \emph{rahasya}, once more
  \item \emph{Say it:} \emph{dāktare ā rognũ rahasya śodhī kāḍhyũ.} — The doctor found out the secret of this illness.
  \item \emph{Recall:} say the Gujarati for an observation, then the Gujarati for proven, then say \emph{rahasya} again
\end{itemize}

\begin{tcolorbox}[breakable,colback=teal!4,colframe=teal!35!black,title={Before you move on}]
What does \gu{રહસ્ય} mean? (\textquotedblleft{}A mystery; a secret\textquotedblright{}.) Say it once more.
\end{tcolorbox}
\section[\texorpdfstring{\gu{ઉઘાડવું} (ughāḍvũ)}{ઉઘાડવું (ughāḍvũ)}]{\gu{ઉઘાડવું} (ughāḍvũ) — to uncover; to open}

\label{lesson:GU-C364-ughadvu-uncover}



\begin{quote}
Before the new one: say the Gujarati for proven, then the Gujarati for a mystery.
\end{quote}

\subsection*{You'll want to know: \gu{ઉઘાડવું}}
\textbf{\gu{ઉઘાડવું}} — \emph{ughāḍvũ} — \textquotedblleft{}to uncover; to open\textquotedblright{}.

\begin{quote}
\gu{મીરાએ} \gu{સવારે} \gu{બારી} \gu{ઉઘાડી}.
\end{quote}

\begin{quote}
\emph{Mīrāe savāre bārī ughāḍī.}
\end{quote}

\begin{quote}
Mira opened the window in the morning.
\end{quote}

Past, with \gu{તેણે} (he or she …): \gu{તેણે} \gu{બારણું} \gu{ઉઘાડ્યું}, \emph{teṇe bārṇũ ughāḍyũ}.

\begin{cousinweb}
\textbf{\gu{ઉઘાડવું}} (\emph{ughāḍvũ}) is the transitive partner of \textbf{\gu{ઊઘડવું}} (\emph{ūghaḍvũ}), \textquotedblleft{}to open up\textquotedblright{}; CDIAL lists it under Sanskrit \emph{udghāṭayati}, \textquotedblleft{}opens\textquotedblright{}, beside a reconstructed \emph{udghaṭati}, \textquotedblleft{}is open\textquotedblright{} (CDIAL 1968).
\end{cousinweb}

\begin{grammarlens}[title={What you've built}]
One more word for giving reasons and opinions.
\end{grammarlens}

\subsection*{Your turn}
\begin{itemize}
\raggedright
  \item \emph{Say it:} \emph{ughāḍvũ}
  \item \emph{Say it:} \emph{ughāḍvũ}, once more
  \item \emph{Say it:} \emph{Mīrāe savāre bārī ughāḍī.} — Mira opened the window in the morning.
  \item \emph{Recall:} say the Gujarati for proven, then the Gujarati for a mystery, then say \emph{ughāḍvũ} again
\end{itemize}

\begin{tcolorbox}[breakable,colback=teal!4,colframe=teal!35!black,title={Before you move on}]
What does \gu{ઉઘાડવું} mean? (\textquotedblleft{}To uncover; to open\textquotedblright{}.) Say it once more.
\end{tcolorbox}
